% MAIN SNAPSHOT — LSPV/RMV::0658::2026-10-02T06:58:00-07:00
% !TEX program = xelatex
\documentclass[11pt]{article}\usepackage[margin=.86in]{geometry}\usepackage{amsmath,amssymb}
\begin{document}
\section*{LSPV-5 MAIN Snapshot 0658}
\textbf{STATUS: OPEN.} P versus NP remains open. No unrestricted SAT lower bound, arbitrary-algorithm transfer theorem, nonrelativizing hinge, or nonalgebrizing hinge is established.

\section*{CORPUS RECONCILIATION}
\textbf{Corpus surveyed.} Governance AGENTS/CODEX/pnp-research skill, root main.tex, MAIN 0502 and materially relevant earlier MAINs, TASK-001--007, specialist chain through SLICECOVER1/MATCIRC1/FACTOR1/TRACECLOSE1, active HEC/VSC-HEC/G-PPC-HEC/CRS-PPC/CCSPC/RCCN/Search--G materials, COVERBREAK1/SOURCEDIR1 experiment artifacts, and post-0502 GROK distills through d3e92831 were reconciled. Repository history shows no newer verified LSPV-1--4 theorem after TRACECLOSE1.

\textbf{New verified results.} No new LSPV theorem since MAIN 0502. TRACECLOSE1 remains PROVED/EXACT: every polynomial-time functional of the ordinary polynomial transcript of a deterministic polynomial-time solver is itself polynomial-time computable. Post-0502 GROK contributes literature classification only.

\textbf{Conflicts resolved.} Root main.tex still points to MAIN 0414 and its LSPV-4 region omits TRACECLOSE1. This snapshot adopts MAIN 0502's correction. COVERBREAK1 remains a COMPUTATIONAL counterexample to universal nonconstant-leaf maximal rank; SOURCEDIR1 remains COMPUTATIONAL finite evidence, not an asymptotic theorem.

\textbf{Stale claims removed/downgraded.} Full polynomial transcript access is not a nonrelativizing hinge. Coding-theoretic shortening/filtration and matroid-covering results remain diagnostics unless a source-matched theorem is proved.

\textbf{Falsifications retained.} Generic short-source/canonicity/locality implications, generic canonical binding, ordinary switching as universal transfer, free atomic shortcuts, optional failed-shortcut costs, raw incidence/tape geometry, generic STORE-or-REACQUIRE, static reread hardness, transcript provenance alone, and SGLAUD1 constant-query amplification remain FALSIFIED/RETIRED.

\textbf{Open arrows changed.} None closed. The acquisition test is sharper: any candidate C(A,x) polynomial-time computable from (x,Tr_A(x)) is retired by TRACECLOSE1; any FP-regenerable selector is retired by REACQ1. A survivor needs an explicit extra primitive plus an implementation-invariant necessity theorem.

\textbf{Adopted/rejected handoffs.} Adopt TRACECLOSE1 and the post-0502 acquisition taxonomy. Adopt recent coding/matroid literature only as ANALOGY/CANDIDATE diagnostics. Reject automatic transfer of restricted proof-system, circuit, matroid, or coding lower bounds to unrestricted SAT runtime.

\section*{Current progress}
The exact structural falsification inequalities remain
\[
\rho(H)\ge1+R+d,\qquad
\rho(H)\ge1+\frac12\left(R+\frac{R^2}{\tau_L}\right).
\]
COVERBREAK1 realizes the finite hostile regime R=m,d=m,tau_L=ceil(m/2) for tested m=5,...,9. SOURCEDIR1 found d=binom(m,2), tau_L<=2 throughout its tested h>m battery, including 822 exhaustive boundary leaves. Therefore the useful source theorem is not maximal rank but the weaker
\[
\boxed{\dim U=\omega(m)}
\]
under the exact reconstruction-compatible vMS parameter regime, or an explicit classification of all low-exterior exceptions.

\section*{Reconstruction, acquisition, and causal state}
FACTOR1 removes YES/NO double-exponent accounting but does not eliminate the dependency exists sigma forall u exists z forall y R. QLEN1 is exact only in the explicit-query-writing model. EXPINV1 and conservative Viola/SGLHYB1 ledgers remain in force. Acquisition and normalization charges remain explicit. TRACECLOSE1 now filters every proposed nonblack-box resource. HEC/CRS-PPC/CCSPC/RCCN remain restricted-model theorems; REREAD1, PCLOSE1-R1, RANDCLOSE1, REACQ1, PADINV1, ISOCAN1, WEIGHTPAD1 and TRACECLOSE1 block current generic promotion to static SAT hardness.

\section*{Barrier and literature audit}
Recent algebrization-barrier work strengthens, rather than removes, the requirement for a genuinely nonalgebrizing hinge. Recent near-maximum circuit lower bounds with adaptive MA/NP query structure show that rounds and witness lengths can be quantitatively meaningful in restricted exponential-time circuit settings, but supply no unrestricted SAT-runtime transfer. Weak-rank proof-system lower bounds and depth-two threshold lower bounds are likewise model-specific. No reviewed primary source closes TC4 or the final P-vs-NP hinge.

\section*{Shortest surviving proof chain}
\[
P=NP\to A\in TIME(N^K)\to H_{SAME}
\to\boxed{\text{exact reconstruction-compatible SU/vMS realization}}
\to\boxed{\text{source-specific }\dim U=\omega(m)\text{ or classified exceptions}}
\]
\[
\to\boxed{\text{charged acquisition/normalization surviving REACQ1/TRACECLOSE1}}
\to\boxed{\text{model-matched reconstruction}}
\to\boxed{\text{arbitrary-algorithm nonrelativizing/nonalgebrizing hinge}}
\to\boxed{\text{unconditional contradiction}}.
\]
Every boxed implication remains OPEN or CONDITIONAL.

\section*{ONE SINGLE TEAM TARGET}
\[
\boxed{\textbf{Source-transcribe one exact reconstruction-compatible SU/vMS line and prove or falsify the three-parameter shortening/innovation statement }\dim U=\omega(m).}
\]
Run COVERBREAK1 as the mandatory hostile regression and SOURCEDIR1 as finite source-direction regression. If structural growth survives, write the resulting acquisition object explicitly as C(A,x) and immediately apply REACQ1/TRACECLOSE1 before attempting any hardness promotion.

\end{document}
% END_OF_VERIFIED_MAIN::0658